\chapter{Ports and Adapters}

\section*{Tujuan Pembelajaran}

\begin{enumerate}
  \item Menjelaskan peran port dan adapter serta arah ketergantungan yang
    membuat logika inti bebas dari teknologi penyimpanan.
  \item Membedakan port masukan dari port keluaran, lalu menentukan komponen
    mana yang menjadi adapter pada sebuah sistem.
  \item Mengukur kemandirian logika inti, biaya penggantian adapter, dan
    cakupan pengujian tanpa infrastruktur.
\end{enumerate}

\section{Pendahuluan}

Bab~2 menunjukkan bahwa \textit{layered architecture} menjaga arah panggilan
tetap ke bawah. Aturan tersebut menahan perubahan di satu lapisan, tetapi tidak
menghilangkan satu persoalan. Lapisan bisnis tetap memanggil lapisan
persistensi, sehingga lapisan bisnis tetap mengenal keberadaan penyimpanan.

Ports and adapters, yang juga disebut \textit{hexagonal architecture},
menyelesaikan persoalan tersebut dengan membalik arah ketergantungan. Logika
inti tidak lagi memanggil penyimpanan. Logika inti justru menyatakan apa yang
dibutuhkannya dalam bentuk antarmuka, lalu pihak lain yang menyediakannya.
Antarmuka itu disebut port, dan pengisinya disebut adapter.

Pembalikan tersebut menghasilkan satu klaim yang tidak biasa, sebab klaim itu
menyebut angka yang pasti. Jumlah teknologi penyimpanan yang disebut oleh logika
inti harus nol. Klaim yang menyebut angka dapat dibuktikan salah, dan bab ini
membuktikannya lewat pengukuran. Pembahasan mengikuti kerangka yang disusun oleh
Richards dan Ford~\cite{richards2020fundamentals}.

\section{Port dan Adapter}

Gambar~\ref{fig:port-adapter} menunjukkan susunan dasarnya. Logika inti berada
di tengah. Seluruh panah ketergantungan menunjuk ke dalam, sehingga tidak ada
panah yang keluar dari logika inti.

\begin{figure}[htbp]
  \centering
  \begin{tikzpicture}[
    kotak/.style={draw, rounded corners=2pt, minimum width=2.1cm,
      minimum height=0.75cm, align=center, font=\scriptsize},
    inti/.style={draw, rounded corners=3pt, minimum width=3.2cm,
      minimum height=1.5cm, align=center, font=\scriptsize, fill=praditagreen!15},
    panah/.style={-{Stealth[length=2mm]}, thick}
  ]
    \node[inti] (inti) {\textbf{Logika Inti}\\
      {\tiny\mdseries menyatakan port}\\
      {\tiny\mdseries yang dibutuhkannya}};

    \node[kotak, fill=blue!8, left=2.2cm of inti] (masuk) {Adapter Masukan\\
      {\tiny\mdseries antarmuka, API}};
    \node[kotak, fill=green!10, above right=0.35cm and 2.2cm of inti] (memori) {Adapter Memori};
    \node[kotak, fill=green!10, below right=0.35cm and 2.2cm of inti] (sqlite) {Adapter SQLite};

    \draw[panah] (masuk) -- node[above, font=\tiny] {port masuk} (inti);
    \draw[panah] (memori.west) -- node[above, font=\tiny, pos=0.45] {port keluar} (inti.east);
    \draw[panah] (sqlite.west) -- (inti.east);
  \end{tikzpicture}
  \caption{Seluruh panah menunjuk ke dalam, sehingga logika inti tidak pernah
    mengenal adapter mana pun}
  \label{fig:port-adapter}
\end{figure}

\textbf{Port} adalah antarmuka yang ditulis oleh logika inti untuk menyatakan
kebutuhannya. Port tidak menyebut teknologi apa pun. Port hanya menyebut
kelakuan yang dibutuhkan, misalnya kemampuan mencari nilai kurs.

\textbf{Adapter} adalah pengisi port tersebut untuk satu teknologi tertentu.
Satu port dapat diisi oleh banyak adapter sekaligus, misalnya satu untuk
penyimpanan dalam memori dan satu lagi untuk basis data.

Port dibedakan menjadi dua arah. \textbf{Port masukan} dipakai pihak luar untuk
memanggil logika inti, misalnya antarmuka pengguna atau API. \textbf{Port
keluaran} dipakai logika inti untuk meminta sesuatu dari luar, misalnya membaca
data. Bab ini menitikberatkan port keluaran, sebab di situlah ketergantungan
pada teknologi paling sering bocor.

Listing~\ref{lst:port} memperlihatkan port keluaran pada contoh bab ini. Port
tersebut hanya menyebut satu metode, dan tidak menyebut basis data sama sekali.

\begin{lstlisting}[style=PythonStyle, caption={Port keluaran yang ditulis oleh
  logika inti, tersedia di \berkas{sources/chapter04/domain.py}}, label=lst:port]
class SumberKurs(Protocol):
  """Port keluaran yang dibutuhkan logika inti.

  Siapa pun yang menyediakan metode ini dapat dipasang, baik penyimpanan dalam
  memori, basis data, maupun tiruan untuk pengujian.
  """

  def nilai(self, kode_asal, kode_tujuan):
    """Mengembalikan nilai kurs, atau None bila pasangannya tidak ada."""
\end{lstlisting}

\section{Menerima Adapter sebagai Argumen}

Port saja belum cukup. Logika inti juga tidak boleh membuat adapternya sendiri.
Bila logika inti memanggil konstruktor sebuah adapter, logika inti kembali
mengenal teknologinya, dan seluruh manfaatnya hilang.

Analogi yang mendekati adalah stopkontak listrik di dinding. Stopkontak
menentukan bentuk colokan yang diterimanya, tetapi tidak menentukan alat apa
yang dicolokkan. Kipas angin, lampu, dan pengisi daya telepon sama-sama dapat
dipasang. Dinding tidak perlu diubah setiap kali alatnya berganti.

Karena itu adapter diserahkan dari luar sebagai argumen.
Listing~\ref{lst:terima-adapter} memperlihatkan caranya. Parameter pertama
adalah sumber kurs, dan fungsi tersebut tidak pernah tahu adapter mana yang
sedang dipakainya.

\begin{lstlisting}[style=PythonStyle, caption={Adapter diterima sebagai argumen,
  tersedia di \berkas{sources/chapter04/domain.py}}, label=lst:terima-adapter]
def hitung_konversi(sumber, kode_asal, kode_tujuan, nominal):
  """Menghitung hasil konversi setelah memeriksa dua aturan bisnis.

  Sumber kurs diterima sebagai argumen, bukan dibuat di dalam fungsi. Karena itu
  fungsi ini tidak pernah mengetahui teknologi penyimpanan yang dipakai.
  """
  if nominal <= 0:
    raise ValueError("Nominal harus lebih besar dari nol")
  if nominal > BATAS_NOMINAL_WAJAR:
    raise ValueError("Nominal melewati batas wajar")
  kurs = sumber.nilai(kode_asal, kode_tujuan)
  if kurs is None:
    raise LookupError(f"Kurs {kode_asal} ke {kode_tujuan} tidak tersedia")
  return nominal * kurs
\end{lstlisting}

Penentuan adapter dipindahkan ke satu berkas perangkai. Pada contoh bab ini
berkas tersebut adalah \berkas{sources/chapter04/aplikasi.py}. Hanya berkas
itulah yang mengetahui kedua adapter sekaligus.

\section{Tiga Akibat yang Dapat Diukur}

Susunan tersebut membawa tiga akibat, dan ketiganya menghasilkan angka.
Tabel~\ref{tab:akibat-port} merangkumnya beserta hasil pengukuran nyata pada
contoh bab ini.

\begin{table}[htbp]
  \centering
  \caption{Tiga akibat susunan port dan adapter beserta hasil pengukurannya}
  \label{tab:akibat-port}
  \begin{tabularx}{\textwidth}{|X|l|X|}
    \hline
    \textbf{Akibat} & \textbf{Angka} & \textbf{Cara mengukur} \\
    \hline
    Logika inti tidak mengenal teknologi & 0 adapter & Menghitung modul yang diimpor \berkas{domain.py} \\
    \hline
    Penggantian adapter tidak menyentuh logika inti & 1 baris & Mengganti pilihan adapter pada berkas perangkai \\
    \hline
    Pengujian berjalan tanpa infrastruktur & 0 basis data & Memasang adapter tiruan berisi satu metode \\
    \hline
  \end{tabularx}
\end{table}

Akibat pertama adalah yang paling penting, sebab angkanya dapat membuktikan
klaim tersebut salah. Bila \berkas{domain.py} ternyata mengimpor
\texttt{sqlite3}, klaim kemandirian gugur pada penerapan tersebut, seberapa pun
rapi susunan berkasnya. Listing~\ref{lst:periksa-port} memperlihatkan
pemeriksaannya.

\begin{lstlisting}[style=PythonStyle, caption={Pemeriksaan kemandirian logika
  inti, tersedia di \berkas{sources/chapter04/periksa_port.py}},
  label=lst:periksa-port]
def modul_diimpor(nama_berkas):
  """Mengumpulkan seluruh nama modul yang diimpor sebuah berkas."""
  with open(nama_berkas, encoding="utf-8") as berkas:
    pohon = ast.parse(berkas.read(), filename=nama_berkas)
  nama = []
  for simpul in ast.walk(pohon):
    if isinstance(simpul, ast.Import):
      nama.extend(a.name for a in simpul.names)
    elif isinstance(simpul, ast.ImportFrom) and simpul.module:
      nama.append(simpul.module)
  return nama
\end{lstlisting}

\section{Kelebihan dan Kekurangan}

Tabel~\ref{tab:untung-rugi-port} merangkum untung rugi susunan ini dibandingkan
\textit{layered architecture} pada Bab~2.

\begin{table}[htbp]
  \centering
  \caption{Perbandingan ports and adapters dengan layered architecture}
  \label{tab:untung-rugi-port}
  \begin{tabularx}{\textwidth}{|l|X|X|}
    \hline
    \textbf{Aspek} & \textbf{Layered} & \textbf{Ports and adapters} \\
    \hline
    Arah ketergantungan & Ke bawah, menuju penyimpanan & Ke dalam, menuju logika inti \\
    \hline
    Pengetahuan logika inti & Mengenal lapisan persistensi & Tidak mengenal teknologi apa pun \\
    \hline
    Penggantian penyimpanan & Menyentuh lapisan persistensi & Menambah satu adapter baru \\
    \hline
    Pengujian & Membutuhkan basis data atau tiruannya & Cukup adapter tiruan berisi satu metode \\
    \hline
    Jumlah berkas & Lebih sedikit & Bertambah oleh port dan adapter \\
    \hline
  \end{tabularx}
\end{table}

Kekurangannya terletak pada jumlah berkas. Setiap kebutuhan keluar menambah satu
port beserta sekurang-kurangnya satu adapter. Untuk aplikasi kecil dengan satu
teknologi penyimpanan yang tidak akan berganti, tambahan tersebut belum tentu
sepadan.

\section{Ringkasan}

Ports and adapters membalik arah ketergantungan yang berlaku pada
\textit{layered architecture}. Logika inti tidak lagi memanggil penyimpanan.
Logika inti menyatakan kebutuhannya sebagai port, lalu adapter mengisinya.
Seluruh panah ketergantungan karena itu menunjuk ke dalam.

Port hanya menyebut kelakuan yang dibutuhkan, bukan teknologi yang
menyediakannya. Satu port dapat diisi banyak adapter sekaligus. Pada contoh bab
ini, adapter memori dan adapter SQLite mengisi port yang sama, dan keduanya
menghasilkan jawaban yang sama tanpa mengubah satu baris pun pada logika inti.

Susunan tersebut hanya bekerja bila logika inti tidak membuat adapternya
sendiri. Adapter diserahkan dari luar sebagai argumen, dan pemilihannya
dipusatkan pada satu berkas perangkai. Berkas perangkai itulah satu-satunya
tempat yang mengetahui seluruh adapter.

Klaim utama pola ini menyebut angka yang pasti, yaitu nol teknologi penyimpanan
yang dikenal logika inti. Klaim yang menyebut angka dapat dibuktikan salah oleh
pemeriksaan sederhana. Pemeriksaan pada bab ini menunjukkan \berkas{domain.py}
hanya mengimpor satu modul standar, sehingga klaimnya terpenuhi.

Keuntungan tersebut ada harganya, yaitu jumlah berkas yang bertambah. Setiap
kebutuhan keluar menambah satu port beserta adapternya. Keputusan memakai pola
ini diambil dengan membandingkan tambahan berkas tersebut terhadap kemungkinan
pergantian teknologi di kemudian hari.

\section{Latihan}

Ketiga latihan berikut memakai satu basis kode yang sama di
\berkas{sources/chapter04/}, tetapi menyasar tiga masalah yang berbeda.
Latihan~1 memeriksa kemandirian logika inti. Latihan~2 mengukur biaya
penggantian adapter. Latihan~3 memeriksa pengujian tanpa infrastruktur.
Penilaian didasarkan pada kelengkapan tabel hasil dan ketepatan jawaban
analisisnya. Hasil setiap mahasiswa dapat berbeda, dan yang dilaporkan adalah
pola perbandingannya.

\subsection*{Latihan 1: Membuktikan Kemandirian Logika Inti}

Latihan ini menguji klaim utama pola tersebut, yaitu logika inti tidak mengenal
teknologi penyimpanan apa pun. Langkah kerjanya sebagai berikut.

\begin{enumerate}
  \item Jalankan pemeriksaan seperti pada Listing~\ref{lst:jalankan-periksa-port},
    lalu catat baris \texttt{Adapter disebut}.
  \item Buka \berkas{sources/chapter04/domain.py}, lalu pastikan tidak ada nama
    teknologi penyimpanan di dalamnya.
  \item Tambahkan baris \texttt{import sqlite3} pada \berkas{domain.py}, lalu
    jalankan ulang pemeriksaannya dan catat hasilnya.
  \item Batalkan penambahan tersebut, lalu pastikan pemeriksaan kembali bersih.
\end{enumerate}

\begin{lstlisting}[language=bash, caption={Menjalankan pemeriksaan kemandirian},
  label=lst:jalankan-periksa-port]
cd sources/chapter04
python periksa_port.py
\end{lstlisting}

Tabel~\ref{tab:lat1port-contoh} memperlihatkan hasil satu kali pemeriksaan nyata.

\begin{table}[htbp]
  \centering
  \caption{Contoh hasil Latihan 1 yang sudah terisi}
  \label{tab:lat1port-contoh}
  \begin{tabularx}{\textwidth}{|X|l|l|}
    \hline
    \textbf{Keadaan berkas} & \textbf{Adapter disebut} & \textbf{Klaim} \\
    \hline
    Apa adanya & 0 & terpenuhi \\
    \hline
    Setelah \texttt{import sqlite3} ditambahkan & 1 & gagal \\
    \hline
  \end{tabularx}
\end{table}

Isikan hasil pemeriksaan pada Tabel~\ref{tab:lat1port-hasil}.

\begin{table}[htbp]
  \centering
  \caption{Lembar isian Latihan 1}
  \label{tab:lat1port-hasil}
  \begin{tabularx}{\textwidth}{|X|l|l|}
    \hline
    \textbf{Keadaan berkas} & \textbf{Adapter disebut} & \textbf{Klaim} \\
    \hline
    Apa adanya & \dots & \dots \\
    \hline
    Setelah impor ditambahkan & \dots & \dots \\
    \hline
    Modul yang tetap boleh diimpor & \multicolumn{2}{l|}{\dots} \\
    \hline
  \end{tabularx}
\end{table}

Jawab dua pertanyaan berikut. Pertama, \berkas{domain.py} tetap mengimpor satu
modul standar, sehingga jelaskan mengapa impor tersebut tidak melanggar klaim.
Kedua, jelaskan mengapa klaim yang menyebut angka lebih berguna daripada klaim
yang hanya menyebut sifat.

\subsection*{Latihan 2: Mengukur Biaya Penggantian Adapter}

Latihan ini membuktikan berapa banyak kode yang berubah ketika teknologi
penyimpanan diganti. Langkah kerjanya sebagai berikut.

\begin{enumerate}
  \item Jalankan aplikasi dengan kedua adapter seperti pada
    Listing~\ref{lst:jalankan-dua-adapter}, lalu pastikan hasilnya sama.
  \item Catat berkas mana saja yang harus dibuka untuk berpindah adapter.
  \item Tulis satu adapter ketiga yang membaca kurs dari berkas teks biasa.
  \item Catat jumlah berkas dan baris yang berubah agar adapter ketiga tersebut
    dapat dipakai.
\end{enumerate}

\begin{lstlisting}[language=bash, caption={Menjalankan aplikasi dengan dua adapter},
  label=lst:jalankan-dua-adapter]
cd sources/chapter04
python aplikasi.py memori
python aplikasi.py sqlite
\end{lstlisting}

Tabel~\ref{tab:lat2port-contoh} memperlihatkan hasil satu kali pengukuran nyata.

\begin{table}[htbp]
  \centering
  \caption{Contoh hasil Latihan 2 yang sudah terisi}
  \label{tab:lat2port-contoh}
  \begin{tabularx}{\textwidth}{|X|l|l|}
    \hline
    \textbf{Kegiatan} & \textbf{Berkas berubah} & \textbf{Baris} \\
    \hline
    Berpindah dari memori ke SQLite & 0 & 0 \\
    \hline
    Menambah adapter ketiga & 2 & 22 \\
    \hline
  \end{tabularx}
\end{table}

Isikan hasil pengukuran pada Tabel~\ref{tab:lat2port-hasil}.

\begin{table}[htbp]
  \centering
  \caption{Lembar isian Latihan 2}
  \label{tab:lat2port-hasil}
  \begin{tabularx}{\textwidth}{|X|l|l|}
    \hline
    \textbf{Kegiatan} & \textbf{Berkas berubah} & \textbf{Baris} \\
    \hline
    Berpindah antar adapter yang sudah ada & \dots & \dots \\
    \hline
    Menambah adapter ketiga & \dots & \dots \\
    \hline
    Berkas \berkas{domain.py} ikut berubah & \multicolumn{2}{l|}{\dots} \\
    \hline
  \end{tabularx}
\end{table}

Jawab dua pertanyaan berikut dengan menunjuk angka pada tabel. Pertama, apakah
\berkas{domain.py} ikut berubah ketika adapter ketiga ditambahkan. Kedua,
bandingkan angka tersebut dengan biaya penggantian basis data pada Bab~2, lalu
jelaskan penyebab perbedaannya.

\subsection*{Latihan 3: Menguji Tanpa Infrastruktur}

Latihan ini membuktikan bahwa logika inti dapat diuji tanpa basis data.
Langkah kerjanya sebagai berikut.

\begin{enumerate}
  \item Jalankan \texttt{python uji\_tanpa\_basis\_data.py}, lalu catat hasil
    ketiga pengujiannya.
  \item Buka \berkas{sources/chapter04/uji_tanpa_basis_data.py}, lalu hitung
    jumlah baris kelas \texttt{SumberTiruan}.
  \item Tambahkan satu pengujian baru untuk nominal yang melewati batas wajar.
  \item Hitung berapa banyak infrastruktur yang disentuh seluruh pengujian
    tersebut.
\end{enumerate}

Tabel~\ref{tab:lat3port-contoh} memperlihatkan hasil satu kali pengujian nyata.

\begin{table}[htbp]
  \centering
  \caption{Contoh hasil Latihan 3 yang sudah terisi}
  \label{tab:lat3port-contoh}
  \begin{tabularx}{\textwidth}{|X|l|l|}
    \hline
    \textbf{Pengujian} & \textbf{Lulus} & \textbf{Basis data disentuh} \\
    \hline
    Hasil benar & ya & 0 \\
    \hline
    Nominal nol ditolak & ya & 0 \\
    \hline
    Kurs hilang ditolak & ya & 0 \\
    \hline
  \end{tabularx}
\end{table}

Isikan hasil pengujian pada Tabel~\ref{tab:lat3port-hasil}.

\begin{table}[htbp]
  \centering
  \caption{Lembar isian Latihan 3}
  \label{tab:lat3port-hasil}
  \begin{tabularx}{\textwidth}{|X|l|l|}
    \hline
    \textbf{Pengujian} & \textbf{Lulus} & \textbf{Basis data disentuh} \\
    \hline
    Hasil benar & \dots & \dots \\
    \hline
    Nominal nol ditolak & \dots & \dots \\
    \hline
    Kurs hilang ditolak & \dots & \dots \\
    \hline
    Nominal melewati batas & \dots & \dots \\
    \hline
    Jumlah baris adapter tiruan & \multicolumn{2}{l|}{\dots} \\
    \hline
  \end{tabularx}
\end{table}

Jawab tiga pertanyaan berikut dengan menunjuk angka pada tabel. Pertama, berapa
baris kode yang dibutuhkan untuk menggantikan seluruh basis data pada
pengujian. Kedua, jelaskan mengapa jumlah baris tersebut menjadi sedikit.
Ketiga, dengan menggabungkan hasil ketiga latihan, tentukan apakah tambahan
port dan adapter sepadan untuk aplikasi sekecil ini, lalu sebutkan angka yang
menjadi dasar keputusannya.
